\section{Preliminaries and Background}
\label{sec:prelims}

We fix a finite set $\vars$ of program variables and a countable state space $\Sigma$; a state $\sigma\in\Sigma$ maps variables to rational values. Arithmetic and Boolean expressions are interpreted as functions on states. For a Boolean expression $b$, let $\llbracket b\rrbracket=\{\sigma\mid \sigma\models b\}$ and let $[b]$ denote its Iverson bracket. Updating $x$ in $\sigma$ to $a(\sigma)$ is written $\sigma[x/a(\sigma)]$.

\subsection{Probabilistic Programs and Subdistributions}

We consider the probabilistic guarded-command language
\[
 \begin{aligned}
 C &::= \ski \mid \ass{x}{a} \mid \comp{C}{C} \mid \pchoice{C}{p}{C}\\
   &\phantom{{}::={}}\mid \ifelse{b}{C}{C} \mid \while{b}{C},
 \end{aligned}
\]
where $p\in[0,1]\cap\mathbb{Q}$. The language includes probabilistic choice but excludes nondeterministic choice. A \emph{subdistribution} is a function $\mu:\Sigma\to\mathbb{R}_{\geq0}$ of weight $\weight{\mu}=\sum_{\sigma}\mu(\sigma)\leq1$. We write $\subdists$ for all subdistributions, $\supp(\mu)=\{\sigma\mid\mu(\sigma)>0\}$ for the support, and $\dirac{\sigma}$ for the Dirac distribution at $\sigma$. Missing mass represents nontermination.

\subsection{Distribution-Transformer Semantics and AST}

Programs denote linear, $\omega$-continuous post-distribution transformers $\post{C}:\subdists\to\subdists$ that do not increase mass. This is the forward denotational semantics of the probabilistic language: each input subdistribution has one output subdistribution, in the tradition of Kozen~\cite{DBLP:conf/focs/Kozen79,DBLP:conf/stoc/Kozen83}. The clauses needed below are
\begin{align*}
 \post{\ski}(\mu) &= \mu,\\
 \post{\ass{x}{a}}(\mu)(\sigma)
   &= \sum_{\sigma':\,\sigma'[x/a(\sigma')]=\sigma}\mu(\sigma'),\\
 \post{\pchoice{C_1}{p}{C_2}}(\mu)
   &= \post{C_1}(p\mu)+\post{C_2}((1-p)\mu),\\
 \post{\comp{C_1}{C_2}}(\mu)
   &= \post{C_2}(\post{C_1}(\mu)),\\
 \post{\ifelse{b}{C_1}{C_2}}(\mu)
   &= \post{C_1}([b]\mu)+\post{C_2}([\neg b]\mu).
\end{align*}
For a loop $C=\while{b}{B}$, we take a least fixed point over transformers, as in~\cite{zilken2026verifyingsamplingalgorithmsdistributional}. We order subdistributions and transformers pointwise. Let $\mathcal{H}$ be the pointed $\omega$-cpo of $\omega$-continuous transformers $f:\subdists\to\subdists$ satisfying $\weight{f(\eta)}\leq\weight{\eta}$ for every $\eta\in\subdists$, with bottom $\bot=\lambda\eta.0$. Define
\begin{equation}
\label{eq:while-transformer-fixpoint}
 \begin{aligned}
  \Phi_{b,B}(f)(\eta)
    &= [\neg b]\eta+f\bigl(\post{B}([b]\eta)\bigr),\\
  \post{C}(\mu)
    &= (\lfp\Phi_{b,B})(\mu)
     = \sup_{j\in\mathbb{N}}\Phi_{b,B}^{j}(\bot)(\mu).
 \end{aligned}
\end{equation}
The functional $\Phi_{b,B}:\mathcal{H}\to\mathcal{H}$ is $\omega$-continuous, and every approximant has mass at most $\weight{\mu}$ (see \Cref{app:translation}). For $j\geq1$, $\Phi_{b,B}^{j}(\bot)(\mu)$ collects the terminating executions with fewer than $j$ completed body iterations. Thus the fixed point ranges over transformers, while all their inputs and outputs remain subdistributions.

Linearity and $\omega$-continuity give
$\post{C}(\sum_i p_i\dirac{\sigma_i})=\sum_i p_i\post{C}(\dirac{\sigma_i})$ for any finite or countable family with $p_i\geq0$ and $\sum_i p_i\leq1$, which lets us reduce one-iteration obligations to Dirac inputs.

\begin{definition}[Almost-sure termination]
A program $C$ is \emph{almost-surely terminating} on $M\subseteq\subdists$ if
$\weight{\post{C}(\mu)}=\weight{\mu}$ for every $\mu\in M$.
\end{definition}

Thus AST means that execution loses no probability mass. We do not require positive AST: the expected runtime may be infinite.

\subsection{Distributional Invariants}

Our assertions are sets of distributions rather than predicates over individual states. This aligns the termination argument with existing partial-correctness proofs for sampling algorithms~\cite{zilken2026verifyingsamplingalgorithmsdistributional}.

\begin{definition}[Distributional invariant]
Let $C=\while{b}{B}$ and $M\subseteq\subdists$. A set $I\subseteq\subdists$ is a \emph{distributional invariant} for $C$ and $M$ if
\[
 M\subseteq I
 \quad\text{and}\quad
 [\neg b]\mu+\post{B}([b]\mu)\in I
 \quad\text{for all }\mu\in I.
\]
We abbreviate $\supp(I)=\bigcup_{\mu\in I}\supp(\mu)$.
\end{definition}

The invariant can simultaneously carry state-space bounds and genuinely distributional information such as equal probability of states in a row of a decision tree. The AST rule uses the support information, while retaining the stronger invariant permits direct composition with partial correctness.

\subsection{Supermartingales, Variants, and Sublevel Sets}

A function $v:\supp(I)\to\mathbb{R}_{\geq0}$ is a one-iteration supermartingale for loop body $B$ if, for every enabled state $\sigma$,
\[
  \sum_{\sigma'\in\Sigma}\post{B}(\dirac{\sigma})(\sigma')v(\sigma')\leq v(\sigma).
\]
For $r\in\mathbb{R}_{\geq0}$, its sublevel set is
$v_{\leq r}=\{\sigma\in\supp(I)\mid v(\sigma)\leq r\}$.
A variant $u:\supp(I)\to\mathbb{N}$ need not decrease in expectation or on every outcome. We only require a uniformly positive chance to move to a smaller value. The role of $v$ is to ensure that executions cannot escape all regions in which $u$ has a finite upper bound.

This division of labour matters for rejection samplers. A rejected run may reset $u$ to a larger value; AST follows because progress remains possible with nonzero probability while $v$ controls the size of the relevant search space.

\subsection{AST Certificates for Probabilistic Control-Flow Graphs}
\label{sec:background}

Our program-level rule builds on the sound and relatively complete
pCFG rule of Majumdar and Sathiyanarayana%
~\cite{DBLP:journals/pacmpl/MajumdarS25}. Their framework includes
bounded demonic nondeterminism and requires AST under every scheduler.
We state the specialization to pCFGs without nondeterministic choice,
matching the language introduced above, before relating its
location-dependent certificates to our distribution-transformer
formulation.

A probabilistic control-flow graph (pCFG) has configurations
$(\ell,\sigma)$ consisting of a control location $\ell$ and a program
state $\sigma$. Its locations are partitioned into assignment locations
$L_{\mathsf A}$, conditional locations $L_{\mathsf C}$, and
probabilistic locations $L_{\mathsf P}$. Assignment and conditional
configurations have a unique enabled successor. At a probabilistic
location, an outgoing edge is selected according to its fixed
probability. The distinguished terminal location
$\ell_{\mathsf{out}}$ is absorbing.

For a configuration $\gamma=(\ell,\sigma)$, let
$\Succ(\gamma)$ denote the set of its successor configurations. If
$\ell\in L_{\mathsf P}$, let
$\Pr(\gamma,\gamma')$ denote the probability of moving from $\gamma$
to $\gamma'\in\Succ(\gamma)$. An invariant
$\mathit{Inv}$ is \emph{inductive} for a set $E$ of initial states if
\[
 \{(\ell_{\mathsf{init}},\sigma)\mid \sigma\in E\}
 \subseteq \mathit{Inv}
\]
and $\mathit{Inv}$ is closed under all positive-probability
transitions.

The rule uses two certificates. A non-negative supermartingale $V$
controls the probability of reaching configurations arbitrarily far
from termination. A natural-valued variant $U$ provides local
progress towards termination. For $r\in\mathbb{R}_{\geq 0}$, define
the $V$-sublevel set
\[
 V_{\leq r}
 =
 \bigl\{\gamma\in\mathit{Inv}\mid V(\gamma)\leq r\bigr\}.
\]

\begin{theorem}
  [pCFG AST rule, adapted from%
  ~\cite{DBLP:journals/pacmpl/MajumdarS25}]
 \label{thm:pcfg-rule}
Let $\mathcal{G}$ be a pCFG with initial location
$\ell_{\mathsf{init}}$ and terminal location
$\ell_{\mathsf{out}}$, and let $E$ be a set of initial states.
The pCFG $\mathcal{G}$ is almost-surely terminating from $E$ if there
exist the following objects:
\begin{enumerate}
 \item An inductive invariant $\mathit{Inv}$ containing $(\ell_{\mathsf{init}},\sigma) \text{for every }\sigma\in E.$

 \item A function
$
  V\colon\mathit{Inv}\to\mathbb{R}_{\geq 0}
$
 satisfying:
 \begin{enumerate}
  \item
  \[
   V(\ell,\sigma)=0
   \quad\Longleftrightarrow\quad
   \ell=\ell_{\mathsf{out}}.
  \]

  \item For every
  $\gamma=(\ell,\sigma)\in\mathit{Inv}$ with
  $\ell\in L_{\mathsf A}\cup L_{\mathsf C}$ and unique successor
  $\gamma'\in\Succ(\gamma)$,
  \[
   V(\gamma')\leq V(\gamma).
  \]

  \item For every
  $\gamma=(\ell,\sigma)\in\mathit{Inv}$ with
  $\ell\in L_{\mathsf P}$,
  \[
   \sum_{\gamma'\in\Succ(\gamma)}
   \Pr(\gamma,\gamma')\,V(\gamma')
   \leq V(\gamma).
  \]
 \end{enumerate}

 \item A function
$
  U\colon\mathit{Inv}\to\mathbb{N}
$
 satisfying:
 \begin{enumerate}
  \item
  \[
   U(\ell,\sigma)=0
   \quad\Longleftrightarrow\quad
   \ell=\ell_{\mathsf{out}}.
  \]

  \item For every
  $\gamma=(\ell,\sigma)\in\mathit{Inv}$ with
  $\ell\in L_{\mathsf A}\cup L_{\mathsf C}$ and unique successor
  $\gamma'\in\Succ(\gamma)$,
  \[
   U(\gamma')<U(\gamma).
  \]

  \item For every $r\in\mathbb{R}_{\geq 0}$:
  \begin{enumerate}
   \item the set
$
    \{U(\gamma)\mid \gamma\in V_{\leq r}\}
$
   is bounded; and

   \item there exists $\varepsilon_r>0$ such that, for every
   $\gamma=(\ell,\sigma)\in V_{\leq r}$ with
   $\ell\in L_{\mathsf P}$,
   \[
    \sum_{\substack{\gamma'\in\Succ(\gamma)\\
                    U(\gamma')<U(\gamma)}}
    \Pr(\gamma,\gamma')
    \geq \varepsilon_r.
   \]
  \end{enumerate}
 \end{enumerate}
\end{enumerate}
Moreover, the rule is relatively complete with respect to the
underlying first-order theory used to express the invariant and the
certificates.
\end{theorem}
The soundness argument has two complementary layers. Fix a sublevel
set $V_{\leq r}$. Since $U$ is natural-valued and bounded on this set,
it ranges over a finite interval. At deterministic locations it
strictly decreases, while at probabilistic locations it decreases
with probability at least $\varepsilon_r$. Consequently, a run that
remains inside $V_{\leq r}$ cannot avoid termination with positive
probability.

It remains to account for runs that leave every fixed sublevel set.
The supermartingale conditions ensure that $V$ does not increase in
expectation. Standard supermartingale arguments therefore make
indefinite escape through increasingly large sublevel sets a
probability-zero event. Combining both layers yields almost-sure
termination. Unlike proof rules based on a globally bounded variant,
the sublevel-set formulation also handles null-recurrent programs,
such as the symmetric one-dimensional random walk.

Although semantically powerful, the pCFG rule is syntactically
fine-grained. Suppose that one source-level loop iteration consists
of a guard, two assignments, a probabilistic choice, and a
conditional. Its pCFG translation introduces a separate control
location for each of these operations. A proof must therefore define
$U$ and $V$ at every intermediate location. Location-dependent
constants, scaling factors, and offsets may be needed solely to
ensure the required inequalities along the internal transitions of
one source-level iteration.

Our rule retains the same two-layer termination argument but moves it
to program level. Instead of assigning certificates to intermediate
control locations, it contracts each complete execution of a
loop-free body into one distribution-transformer step. The
certificates can therefore express progress over complete loop
iterations, matching the abstraction at which sampling algorithms
are naturally understood.

\paragraph{Translation assumption.}
The standard structural translation $\mathcal{T}[C]$ maps program
syntax to a pCFG. Assignments become update edges, conditionals become
guarded edges, and probabilistic choices become probabilistic
locations. Sequential composition connects the exit of one component
to the entry of the next, while a while-loop connects the body exit
back to the guard.

For a loop-free body $B$, paths through $\mathcal{T}[B]$ from its
entry location $\ell_{\mathsf{in}}^B$ to its exit location
$\ell_{\mathsf{out}}^B$ correspond exactly to the outcomes of its
distribution-transformer semantics:
\[
 \post{B}(\dirac{\sigma})(\sigma')
 =
 \Pr\bigl(
   \text{paths in }\mathcal{T}[B]
   \text{ from }(\ell_{\mathsf{in}}^B,\sigma)
   \text{ to }(\ell_{\mathsf{out}}^B,\sigma')
 \bigr).
\]
Consequently, a program $C$ is AST on a set $M$ of initial
subdistributions exactly when $\mathcal{T}[C]$ reaches its terminal
location almost surely from every state in $\supp(M)$. A structural
proof of this correspondence is summarized in
\Cref{app:translation}. This correspondence provides the bridge used
in \Cref{sec:metatheory} to transfer soundness and relative
completeness between the pCFG rule and our program-level rule.
